\subsection{任意集合的外测度}

定义 4. --- 设 \(\mu\) 是 X 上的正测度；对 X 的每个子集 A ，上积分 \(\mu^{*}(\varphi_{\mathrm{A}})\) 称为 A （关于测度 \(\mu\) ）的外测度，记作 \(\mu^{*}(\mathrm{A})\) 。

这样，一个集合的外测度是一个数 \(\geqslant0\) ，有限或等于 \(+\infty\) ，且对开集而言，它与 No.~2 的定义 2 中所定义的外测度一致。

命题 16. --- 若 A 与 B 是 X 的两个子集，满足 A ⊂ B，则 \(\mu^{*}(A) \leqslant \mu^{*}(B)\) 。

推论. --- X 中每个相对紧集的外测度都是有限的。

事实上，这样的集合含于一个相对紧的开集（GT, I, §9, No.~7, 命题 10），而后者的外测度是有限的（No.~2, 命题 5）。

命题 17. --- 若 \((\mathrm{A}_{n})\) 是 X 的子集的一个递增序列，则 \(\mu^{*}(\bigcup_{n}\mathrm{A}_{n})=\sup_{n}\mu^{*}(\mathrm{A}_{n})\) 。

命题 18. --- 对 X 的子集的每个序列 (A \(_{n}\) )，

\[
\mu^ {*} \left(\bigcup_ {n} \mathrm{A} _ {n}\right) \leqslant \sum_ {n} \mu^ {*} (\mathrm{A} _ {n}).
\]

这些命题就是 No.~3 的命题 10 与 13 以及定理 3 对集合的特征函数而言的翻译。

命题 19. --- 对 X 的每个子集 A ，\(\mu^{*}(A)\) 是包含 A 的诸开集的外测度的下确界。

若 \(\mu^{*}(\mathbf{A}) = +\infty\) ，命题是显然的。在相反的情形，对每个 \(\varepsilon\) （满足 \(0 < \varepsilon < 1\) ），存在函数 \(f \in \mathcal{I}_{+}\) 使 \(\varphi_{\mathrm{A}} \leqslant f\) 且 \(\mu^{*}(\mathrm{A}) \leqslant \mu^{*}(f) \leqslant \mu^{*}(\mathrm{A}) + \varepsilon\) 。设 \(\mathbf{G}\) 为诸 \(x \in \mathbf{X}\) （满足 \(f(x) > 1 - \varepsilon\) ）所成之集。由于 \(f\) 下半连续，\(\mathbf{G}\) 是开集（GT, IV, §6, No.~2, 命题 1）且包含 A；另一方面 \(f \geqslant (1 - \varepsilon)\varphi_{\mathrm{G}}\) ，由此

\[
\mu^ {*} (\mathrm{G}) \leqslant \frac {1}{1 - \varepsilon} \mu^ {*} (f) \leqslant \frac {1}{1 - \varepsilon} \left(\mu^ {*} (\mathrm{A}) + \varepsilon\right);
\]

由于 \(\varepsilon\) 任意，可见 \(\mu^{*}(\mathbf{G})\) 与 \(\mu^{*}(\mathbf{A})\) 可相差得任意小，由此得命题。

